\scanpage{120}\mspage{65}
On sait d'ailleurs que si un foncteur \inserted{$f$} a un adjoint $g$ (à gauche ou à droite), alors $f$ est un foncteur localisation \unsure{sss|ssi} l'adjoint $g$ est pl. fidèle. Pour prouver que $i,j$ sont pl. fidèles, il serait donc tentant de prouver qu'ils sont adjoints des foncteurs $\unsure{\rho|p},\sigma$.

Pour \deleted{\unsure{voir}} y \unsure{venir à bout|voir un peu plus clair}, je vais poser
\[
 \left\{\begin{aligned}
 \deleted{\unsure{\mathcal A}}\unsure{A|\mathcal A}&=\mathrm{Cat}_{/S},\quad
 \deleted{\unsure{W_0}}\Sigma_0=W_S^d,\quad\Sigma=W_S\qquad(\subset\mathrm{Fl}(A))\\
 C:\ A&\unsure{\longrightarrow|\text{---}}A\quad\text{le foncteur familier, qui a la vertu}\\
 &C(\Sigma)\subset\Sigma_0
 \end{aligned}\right.
\]
On a donc la situation
\diagram{\begin{tikzcd}[ampersand replacement=\&,column sep=large]
 {(\unsure{\alpha|d})}\& (A,\Sigma_0)\arrow[r,shift left,"\rho_0=\mathrm{id}_A"]\&
 (A,\Sigma),\arrow[l,shift left,"C"]
\end{tikzcd}}
induisant par passage aux localisés
\diagram{\begin{tikzcd}[ampersand replacement=\&,column sep=large]
 A\Sigma_0^{-1}\arrow[r,shift left,"\rho=\overline{\rho_0}"]\&
 A\Sigma^{-1},\arrow[l,shift left,"\overline C"]
\end{tikzcd}}
on veut \deleted{\unsure{dégager}} vérifier que $\rho$ et $C$ sont adjoints l'un de l'autre, $C$ \uuline{sûrement à droite}, ce qui signifie que l'on doit avoir des hom fonctoriels
% slanted
\margin{mais si $C$ est à \uline{droite} les flèches ici doivent aller en sens inverse et c'est bien à quoi je dois \unsure{en venir|me résoudre} \unsure{on dirait|on dit} p.~67}
\diagram{\begin{tikzcd}[ampersand replacement=\&,row sep=small,column sep=small]
 {}\arrow[d,phantom,"{\Bigg\{}" description] \&
 \mathrm{id}_{A\Sigma^{-1}}\arrow[r,shift left]\&
 \rho\overline C\arrow[l,shift left,very thick]\& (= \overline{\rho_0 C})\\
 {}\&\mathrm{id}_{A\Sigma_0^{-1}}\arrow[r,shift right,very thick]\&
 \overline C\rho\arrow[l,shift right]\&\unsure{\neq|{}}(=\overline{C\rho_0})
\end{tikzcd}}
satisfaisant deux relations de compatibilité sur lesquelles on reviendra. On peut espérer trouver ces \deleted{\unsure{foncteurs}} hom fonctoriels à partir de
\[
 \left\{\begin{aligned}
 \mathrm{id}_A&\xrightarrow{i}\rho_0 C=C\\
 \mathrm{id}_A&\xleftarrow{\unsure{p|{}}} C\rho_0=C
 \end{aligned}\right.
\]
% archive sheet number: 119
